\chapter{自律走行システムの構成}
\label{chap:system}
本章では，対象とするモビリティとセンサ，および本研究で評価する自律走行システムの構成を述べる．本論文で寄与を評価する要素は，経路生成，経路の平滑化，局所経路計画，速度計画，経路追従制御，停止である．本章では各要素の構成と式を述べ，設定値は3章のTable~\ref{tab:params}にまとめて示す．

\section{モビリティとセンサ}
\label{sec:vehicle}
対象とするモビリティは，AI-Formula運営から提供される三輪モビリティである．2つの前駆動輪はインホイールモータで駆動され，後輪は従動輪（キャスタ）である．トレッドは0.6\,m，ホイールベースは0.733\,mである．車両の寸法は競技規則で全長2000\,mm，全幅800\,mm，全高600\,mm以内と定められている．従動輪の舵角は駆動輪のトルク差によって決まり，直接は操舵できない．従動輪には舵角を観測するエンコーダ（14\,bit，100\,Hz）が設置されている．従動輪には，藤野\cite{fujino2026}による，トレールとばねプリロード，舵角の可動域を制御する機構が備わっており，本研究ではこれらの機構を固定して用いる．

センサは，ステレオカメラ（ZED，640$\times$360画素，15\,Hz，歪み補正済みのカラー画像を使用），GNSS内蔵の慣性航法装置（VectorNav VN-200，IMU，GNSSとも出力20\,Hz），駆動輪の回転速度（CAN経由），従動輪の舵角である．計算機はJetson AGX Orinであり，画像からの推論はTensorRTで実行する．Table~\ref{tab:sensors}に使用するセンサと観測量を示す．自己位置の推定には，駆動輪の回転速度と，慣性航法装置が出力する姿勢を用いる．この姿勢は，装置の内部でGNSSを含む観測から推定されたものである．GNSSの位置は，評価用の参照としてのみ記録する．

\begin{table}[H]
  \centering
  \caption{Sensors and observations used in the system.}
  \label{tab:sensors}
  \begin{tabular}{lll}
    \toprule
    センサ & 観測量 & 用途 \\
    \midrule
    ZEDカメラ & RGB画像 640$\times$360，15\,Hz & 経路生成，信号・障害物の検出 \\
    VN-200 & 角速度，加速度，姿勢，速度 & 自己位置推定（姿勢） \\
    駆動輪エンコーダ & 左右車輪の回転速度 & 速度，自己位置推定 \\
    従動輪エンコーダ & 従動輪の舵角 & 車両応答の同定 \\
    GNSS（VN-200） & 緯度経度 & 評価用の参照位置のみ \\
    \bottomrule
  \end{tabular}
\end{table}

\section{システムの構成}
\label{sec:architecture}
システムの構成をFig.~\ref{fig:system}に示す．システムは，Sensing，Perception，Localization，Planning，Control，Vehicleの6つのモジュールからなる．Sensingはセンサの出力を各モジュールに渡す．Perceptionはカメラ画像から，交差点での進行方向ごとの目標経路と，信号とコーンの位置から決まる停止点を出力する．Localizationは駆動輪の回転速度と慣性航法装置の姿勢から自己位置を求め，運営提供のトポロジカルマップと自己位置から経路選択の指示を生成する．Planningは，目標経路を時間方向に平滑化した参照経路，障害物を回避する局所経路，局所経路に沿った速度の列を順に生成する．Controlは局所経路と速度の列に追従する並進速度と舵角の指令値を出力し，車体制御系で左右駆動輪の回転速度の指令値に変換してVehicleのモータドライバに送る．各処理はROS~2のノードとして実装し，1つのプロセスで実行する．以下の各節で，各モジュールの処理を述べる．

\begin{figure}[H]
  \centering
  \begin{tikzpicture}[
  font=\fontspec{TeX Gyre Heros Cn}\fontsize{8}{8.8}\selectfont,
  >={Stealth[length=1.6mm,width=1.2mm]},
  mod/.style={draw=black!55, line width=0.5pt, fill=black!4, anchor=north west, inner sep=0pt},
  nd/.style={draw=black!85, line width=0.4pt, fill=white, align=center, inner xsep=1.5pt, inner ysep=1.5pt, minimum height=1.0cm},
  lab/.style={font=\fontspec{TeX Gyre Heros Cn}\fontsize{7.5}{8.4}\selectfont, align=center, inner sep=1pt},
  arr/.style={->, line width=0.5pt}
]
  % ---------- 上段: 5 つのモジュール (高さ共通) ----------
  \node[mod, minimum width=1.44cm, minimum height=3.74cm] (sens) at (0,0) {};
  \node[mod, minimum width=1.44cm, minimum height=3.74cm, xshift=0.95cm] (perc) at (sens.north east) {};
  \node[mod, minimum width=4.54cm, minimum height=3.74cm, xshift=1.25cm] (plan) at (perc.north east) {};
  \node[mod, minimum width=2.72cm, minimum height=3.74cm, xshift=1.30cm] (ctrl) at (plan.north east) {};
  \node[mod, minimum width=1.20cm, minimum height=3.74cm, xshift=1.20cm] (veh)  at (ctrl.north east) {};

  \node[anchor=north west, inner sep=0pt] at ($(sens.north west)+(0.07,-0.07)$) {\bfseries Sensing};
  \node[anchor=north west, inner sep=0pt] at ($(perc.north west)+(0.07,-0.07)$) {\bfseries Perception};
  \node[anchor=north west, inner sep=0pt] at ($(plan.north west)+(0.07,-0.07)$) {\bfseries Planning};
  \node[anchor=north west, inner sep=0pt] at ($(ctrl.north west)+(0.07,-0.07)$) {\bfseries Control};
  \node[anchor=north west, inner sep=0pt] at ($(veh.north west)+(0.07,-0.07)$)  {\bfseries Vehicle};

  % Sensing (縦 3 つ)
  \node[nd, minimum width=1.28cm] (cam) at ($(sens.north)+(0,-0.88)$) {Camera};
  \node[nd, minimum width=1.28cm] (enc) at ($(sens.north)+(0,-2.00)$) {Wheel and\\caster\\encoders};
  \node[nd, minimum width=1.28cm] (imu) at ($(sens.north)+(0,-3.12)$) {Inertial\\navigation\\unit};

  % Perception (縦 2 つ)
  \node[nd, minimum width=1.28cm] (pg)  at ($(perc.north)+(0,-0.88)$) {Path\\generation\\(CNN)};
  \node[nd, minimum width=1.28cm] (det) at ($(perc.north)+(0,-2.00)$) {Signal and\\obstacle\\detection};

  % Planning (横 3 つ)
  \node[nd, minimum width=1.24cm] (ps) at ($(plan.north west)+(0.73,-0.88)$) {Path\\smoothing};
  \node[nd, minimum width=1.24cm, right=0.30cm of ps] (lp) {Local path\\planning};
  \node[nd, minimum width=1.24cm, right=0.30cm of lp] (vp) {Velocity\\planning};
  \draw[arr] (ps) -- (lp);
  \draw[arr] (lp) -- (vp);

  % Control (横 2 つ)
  \node[nd, minimum width=1.10cm] (pf) at ($(ctrl.north west)+(0.66,-0.88)$) {Path\\following};
  \node[nd, minimum width=1.10cm, right=0.30cm of pf] (cc) {Chassis\\control};
  \draw[arr] (pf) -- (cc);

  % Vehicle
  \node[nd, minimum width=1.04cm] (iwm) at ($(veh.north)+(0,-0.88)$) {In-wheel\\motors};

  % ---------- 下段: Localization ----------
  \coordinate (L1) at ($(perc.south west)+(0,-1.0)$);
  \coordinate (L2) at ($(ctrl.south east)+(0,-2.35)$);
  \node[mod, inner sep=0pt, fit=(L1)(L2), anchor=center] (loc) {};
  \node[anchor=north west, inner sep=0pt] at ($(loc.north west)+(0.07,-0.07)$) {\bfseries Localization};
  \node[nd, minimum width=1.30cm, minimum height=0.60cm] (odo) at ($(loc.north west)+(1.75,-0.82)$) {Odometry};
  \node[nd, minimum width=1.95cm, minimum height=0.60cm] (route) at (plan.north |- odo) {Route instruction};
  \draw[arr] (odo) -- (route);

  \node[nd, fill=black!10, minimum width=1.95cm, minimum height=0.60cm] (topo) at ($(route.south |- loc.south)+(0,-0.95)$) {Topological map};
  \draw[arr] (topo.north) -- (topo.north |- loc.south);

  % ---------- 上段のモジュール間矢印 (同じ高さで水平) ----------
  \draw[arr] (sens.east |- cam.center) -- node[lab, below=2pt] {camera\\image} (perc.west |- cam.center);
  \draw[arr] (perc.east |- cam.center) -- node[lab, below=2pt] {target path,\\obstacles,\\stop point} (plan.west |- cam.center);
  \draw[arr] (plan.east |- cam.center) -- node[lab, below=2pt] {local path\\with velocity\\profile} (ctrl.west |- cam.center);
  \draw[arr] (ctrl.east |- cam.center) -- node[lab, below=2pt] {wheel\\speed\\commands} (veh.west |- cam.center);

  % ---------- Localization からの上向き矢印 ----------
  \draw[arr] (perc.south |- loc.north) -- node[lab, right=1pt] {route instruction} (perc.south);
  \draw[arr] (plan.south |- loc.north) -- node[lab, right=1pt] {odometry, route instruction} (plan.south);
  \draw[arr] (ctrl.south |- loc.north) -- node[lab, right=1pt] {odometry} (ctrl.south);

  % ---------- Sensing -> Localization ----------
  \draw[arr] (sens.south) -- (sens.south |- odo.center) -- node[lab, above=1pt] {wheel speed,\\attitude} (loc.west |- odo.center);
\end{tikzpicture}

  \caption{Configuration of the autonomous driving system.}
  \label{fig:system}
\end{figure}

\section{カメラ画像に基づく経路生成}
\label{sec:vision_planner}
カメラ画像を直接入力とし，車両座標系の経路を出力する畳み込みニューラルネットワークを用いる．以下，この出力を経路生成の出力（目標経路）と呼び，平滑化後の経路を参照経路，局所経路計画の出力を局所経路と呼ぶ．ネットワークは，画像の縦方向に設けた行ごとに経路が通る横位置を推定する行アンカー形式\cite{qin2020}であり，直進，左折，右折の3つの進行方向に対応する3本の経路と，各経路の有無を同時に出力する．交差点での経路選択の指示は，ネットワークの外で3本のうち1本を選ぶことで与える．出力した画像上の経路点は，カメラの校正値と路面が平面であるという仮定のもとで車両座標系に投影する．

学習データは，実車で収集した画像に対し，GNSSによる姿勢と自作のベクターマップから車線中心を投影して作成した経路を教師とする．GNSSと自作地図について競技規則が禁止するのは走行中の使用であり，オフラインでの教師の作成には用いている．学習には，人手で修正した実画像と，生成モデル（Cosmos-Transfer2.5）による天候の変換で作成した画像を用い，学習用と検証用に分ける．検証の指標が最小となったエポックの重みを採用する．画像の枚数はTable~\ref{tab:params}に示す．

\section{経路の平滑化}
\label{sec:path_smoother}
経路生成の出力はフレームごとに揺らぐため，時間窓内の出力を自己位置推定の変位で現在の車両座標系に揃え，横位置$y$を進行方向$x$の多項式
\begin{equation}
  y(x) = \sum_{k=0}^{n} c_k x^k
  \label{eq:smooth}
\end{equation}
で最小二乗法により当てはめて平滑化する．当てはめに用いる点の範囲（当てはめ範囲）は車両の前方に限る．ただし，経路生成が点を出力するのは最遠の点までの前方距離（以下，認識距離）までであり，当てはめ範囲のうち実際に点があるのはそこまでである．平滑化後の経路は，一定の間隔で前方の出力長まで出力するため，最遠の点より先の部分は多項式を延長した外挿である．平滑化した経路を参照経路として局所経路計画に渡す．

\section{Frenet座標系に基づく局所経路計画}
\label{sec:local_planner}
参照経路に沿ったFrenet座標系\cite{werling2010}で，横変位を5次多項式で変化させる候補経路を生成し，曲率，横変位，障害物との距離からなるコストで1本を選ぶ．障害物がないときは，参照経路からの横変位$d_{\mathrm{lat}}$を，合流距離$L_{\mathrm{m}}$の間に0にする経路を生成する．合流距離は車速$v$に比例させ，合流時の横加速度が上限$a_{y,\mathrm{m}}$を超えないように下限を与え，さらに最小値と最大値の範囲に制限する．
\begin{equation}
  L_{\mathrm{m}} = \mathrm{clip}\left(\max\left(t_{\mathrm{m}} v,\; v\sqrt{\frac{10}{\sqrt{3}}\frac{|d_{\mathrm{lat}}|}{a_{y,\mathrm{m}}}}\right),\; L_{\min},\; L_{\max}\right)
  \label{eq:merge}
\end{equation}
ここで$t_{\mathrm{m}}$は合流時間，第2項は，5次多項式で横変位$d_{\mathrm{lat}}$を距離$L$の間に動かすときの横加速度の最大値$(10/\sqrt{3})v^2|d_{\mathrm{lat}}|/L^2$が$a_{y,\mathrm{m}}$以下となる最小の距離である．障害物があるときは，検出した障害物の左右いずれかを通る経路を生成し，回避する側を通過まで保持する．回避による横変位は車線幅の内側に制限する．経路は毎周期，計測した自己位置から再計算する．局所経路の長さは，一定の長さ（Table~\ref{tab:params}）と参照経路の終端までの長さのうち短い方とする．

\section{曲率と停止点に基づく速度計画}
\label{sec:speed_planner}
局所経路の各点$i$について，Menger曲率$\kappa_i$から横加速度の上限$a_{y,\max}$で決まる速度の上限
\begin{equation}
  v_{\mathrm{curve},i} = \sqrt{a_{y,\max}/\max(|\kappa_i|, \epsilon)}
  \label{eq:vcurve}
\end{equation}
と，停止点までの距離から決まる上限
\begin{align}
  v_{\mathrm{stop},i} &= \sqrt{2 \tilde{a}_{\mathrm{stop}} \max(s_{\mathrm{stop}} - s_i, 0)} \label{eq:vstop}\\
  \tilde{a}_{\mathrm{stop}} &= \min\left(\max\left(a_{\mathrm{stop}},\ \frac{v^2}{2(s_{\mathrm{stop}} - s_0)}\right),\ a_{\mathrm{stop,max}}\right) \label{eq:astop}
\end{align}
を求める．ここで$s_i$は局所経路の始点から点$i$までの弧長，$s_0$は自車位置の弧長，$s_{\mathrm{stop}}$は停止点までの弧長，$v$は自車の速度である．停止点に向けた減速には，走行時の縦加速度の上限$a_{x,\max}$とは別に，停止専用の減速度$a_{\mathrm{stop}}$を用い，その減速度では停止できない場合は必要な減速度$v^2/(2(s_{\mathrm{stop}}-s_0))$まで強め，上限$a_{\mathrm{stop,max}}$で制限する．上限速度の設定値$v_{\mathrm{cap}}$との最小値を各点の上限とする．
\begin{equation}
  v^{\max}_i = \min\left(v_{\mathrm{cap}},\ v_{\mathrm{curve},i},\ v_{\mathrm{stop},i}\right)
  \label{eq:vlimit}
\end{equation}
そのうえで，進行方向に沿った加速側の走査と逆方向の減速側の走査（forward--backward pass）\cite{velenis2008,lipp2014}により，縦加速度の上限$a_{x,\max}$を満たす速度の列を得る．
\begin{align}
  v_{i+1} &\leftarrow \min\left(v_{i+1}, \sqrt{v_i^2 + 2 a_{x,\max}\Delta s_i}\right) \label{eq:forward}\\
  v_{i} &\leftarrow \min\left(v_{i}, \sqrt{v_{i+1}^2 + 2 a_{x,\max}\Delta s_i}\right) \label{eq:backward}
\end{align}
経路の終端では，次の周期で経路が延びることを前提に，終端の曲率から決まる上限をそのまま用いる．停止点は信号認識と障害物検出から与えられる（\ref{sec:perception}節）．交差点の急カーブで横加速度が上限$a_{y,\max}$を超えないように，経路選択の指示が直進以外のあいだは上限速度を下げる（分岐での速度抑制）．速度の列は経路点とともに経路追従制御に渡される．

\section{経路追従制御}
\label{sec:tracker}
経路追従制御は，Pure Pursuit法\cite{coulter1992}と，横方向のモデル予測制御と縦方向のPID制御を組み合わせた制御器（以下，MPCと呼ぶ）を設定で選んで用いる．いずれも並進速度の指令値と舵角の指令値を出力し，並進速度は速度計画の列から一定時間先の値を参照する．舵角の指令値は上限と下限で制限する．

Pure Pursuit法は，注視距離$L_{\mathrm{d}}$の先の注視点を通る円弧の舵角
\begin{equation}
  \delta = k_{\delta} \arctan\frac{2 L \sin\alpha}{L_{\mathrm{d}}},\quad
  L_{\mathrm{d}} = \mathrm{clip}\left(k_{\mathrm{d}} v,\ L_{\mathrm{d,min}},\ L_{\mathrm{d,max}}\right)
  \label{eq:pp}
\end{equation}
を指令する．ここで$L$はホイールベース，$\alpha$は注視点の方向と車両の向きの差，$k_{\mathrm{d}}$は比例係数であり，注視距離は速度に比例させて上下限で制限する．$k_{\delta}$は舵角のゲインであり，円弧の舵角に掛ける．MPCは，操舵応答を時定数$\tau_{\mathrm{mpc}}$の一次遅れ
\begin{equation}
  \dot{\delta}(t) = \frac{\delta_{\mathrm{cmd}}(t) - \delta(t)}{\tau_{\mathrm{mpc}}}
  \label{eq:mpc}
\end{equation}
として内部モデルに持ち，経路に対する横偏差と方位誤差を予測区間で最小化する．

\section{車体制御系}
\label{sec:chassis}
舵角の指令値と並進速度の指令値から，差動二輪モデルで左右駆動輪の回転速度の指令値を生成し，モータドライバに送信する\cite{fujino2026}．

\section{信号と障害物の検出}
\label{sec:perception}
信号パネルとコーンは，カメラ画像に対する物体検出ネットワークで検出し，検出枠の下端を路面に投影して車両座標系の位置を求める．信号が赤と判定されたときは，パネルの前に置かれたコーンの手前に停止点を発行し，青と判定されるまで発行を続ける．コーンは局所経路計画の障害物として扱い，回避できないときは手前に停止点を発行する．停止点は毎フレーム発行し，発行が一定時間（停止点の保持時間，Table~\ref{tab:params}）途切れると，速度計画は停止点なしとして扱う．

\section{自己位置推定と経路選択の指示}
\label{sec:localization}
自己位置は，駆動輪の回転速度から求めた並進速度と慣性航法装置の姿勢を積分するオドメトリ（デッドレコニング）で求め，地図との照合は行わない．本論文では，この処理を自己位置推定と呼ぶ．青コースでは，運営提供のトポロジカルマップと自己位置から，次の交差点での経路選択（直進・左折・右折）の指示を生成し，\ref{sec:vision_planner}節の3本の経路から1本を選ぶ．赤コースでは常に直進を選ぶ．
